\chapter{实验设计与结果分析}
\label{chap:exp}

本章通过系统性的实验验证本文提出的基于驱动程序替换的固件仿真方法的有效性。实验从多个维度对方法进行评估，包括静态分析模块的识别准确率、运行正确性、方法对比以及应用验证。

\section{实验设计}
\label{sec:design}

为了全面评估本文方法的有效性和实用性，本节介绍实验的整体设计，包括实验环境的搭建、测试用例的选择以及实验目标的定义。

\subsection{实验环境}

实验在以下软硬件环境中进行，具体配置如表~\ref{tab:hw_env}和表~\ref{tab:sw_env}所示。

\begin{table}[htbp]
\centering
\caption{硬件环境配置}
\label{tab:hw_env}
\small
\begin{tabular}{ll}
\toprule
\textbf{配置项} & \textbf{规格} \\
\midrule
主机 & Intel Xeon Platinum 8350C 处理器（128核心），503GB 内存，869GB SSD \\
\bottomrule
\end{tabular}
\end{table}

\begin{table}[htbp]
\centering
\caption{软件环境配置}
\label{tab:sw_env}
\small
\begin{tabular}{ll}
\toprule
\textbf{配置项} & \textbf{规格} \\
\midrule
操作系统 & Ubuntu 22.04 LTS \\
编译工具链 & ARM GCC 10.3.1、CMake 3.21.4 \\
仿真环境 & Unicorn 2.0.2、UnicornAFL（基于Unicorn的模糊测试扩展） \\
模糊测试 & AFL++ 4.08c+（基于v4.08c的开发版本） \\
大模型 & DeepSeek V3.2 \\
\bottomrule
\end{tabular}
\end{table}

\subsection{测试用例}

为了验证方法的通用性和有效性，本研究从 GitHub 收集的开源的物联网设备固件项目涵盖网络服务、环境监测、工业控制等多种应用场景，这些固件作为测试案例，涵盖不同的芯片平台、操作系统、驱动复杂度和应用场景。

\textbf{（1）芯片平台选择}

测试固件覆盖了物联网设备中常见的两类主流 MCU 平台：STM32F4 系列（ARM Cortex-M4）和 NXP i.MX RT 系列（ARM Cortex-M7）。

\textbf{（2）操作系统选择}

测试固件覆盖了三种典型的嵌入式软件架构：FreeRTOS、RT-Thread 和裸机。FreeRTOS 是全球使用最广泛的开源实时操作系统，具有内核精简、易于移植的特点，支持任务管理、队列、信号量等功能。RT-Thread 是国产开源实时操作系统，提供完整的设备驱动框架和丰富的软件包生态，在国内物联网领域应用广泛。裸机是指无操作系统支持的固件，所有驱动和应用代码直接运行在硬件上，常见于资源极度受限的物联网设备。

\textbf{（3）外设类型选择}

测试固件涵盖了物联网设备中常见的四类外设驱动，详细信息如表~\ref{tab:peripherals}所示。

\begin{table}[htbp]
\centering
\caption{外设类型分类}
\label{tab:peripherals}
\begin{adjustbox}{max width=\textwidth}
\begin{tabular}{ll}
\toprule
\textbf{类别} & \textbf{外设类型} \\
\midrule
通信外设 & UART（串口通信）、I2C（总线通信）、SPI（串行外设接口）、ETH（以太网控制器） \\
存储外设 & Flash（闪存控制器）、SDIO（SD 卡接口） \\
定时器外设 & TIM（通用定时器）、RTC（实时时钟）、WDT（看门狗定时器） \\
其他外设 & BLE（蓝牙）、SD（SD卡） \\
\bottomrule
\end{tabular}
\end{adjustbox}
\end{table}

\textbf{（4）测试固件构建}

本研究测试固件来源于官方示例与开源项目，从 STM32CubeMX 和 NXP MCUXpresso SDK 提取的官方示例程序代表规范的驱动代码编写风格，而从 GitHub 和 Gitee 收集的开源的物联网设备固件项目涵盖智能照明、环境监测、工业控制等多种应用场景。这种来源选择确保了测试固件的多样性和代表性。

经过筛选和整理，最终构建了包含 19 个固件的测试集，详细信息如表~\ref{tab:firmware}所示。

\begin{table}[htbp]
\centering
\caption{测试固件详细信息}
\label{tab:firmware}
\scriptsize
\begin{tabularx}{\textwidth}{>{\raggedright\arraybackslash}X p{2.0cm} p{1.5cm} >{\raggedright\arraybackslash}p{1.5cm} >{\raggedright\arraybackslash}p{1.5cm}}
\toprule
\textbf{Demo名称} & \textbf{MCU型号} & \textbf{OS} & \textbf{外设} & \textbf{驱动库} \\
\midrule
\url{NXP_FATFS_BareMetal} & RT1052 & 裸机 & SD & NXP \\
\url{NXP_FATFS_FreeRTOS} & RT1052 & FreeRTOS & SD & NXP \\
\url{NXP_I2C_BareMetal} & RT1052 & 裸机 & I2C & NXP \\
\url{NXP_I2C_FreeRTOS} & RT1052 & FreeRTOS & I2C & NXP \\
\url{NXP_LwIP_BareMetal} & RT1052 & 裸机 & ETH & NXP \\
\url{NXP_LwIP_FreeRTOS} & RT1052 & FreeRTOS & ETH & NXP \\
\url{NXP_UART_BareMetal} & RT1052 & 裸机 & UART & NXP \\
\url{NXP_UART_FreeRTOS} & RT1052 & FreeRTOS & UART & NXP \\
\url{imxrt1052-nxp-evk} & RT1052 & RT-Thread & 多外设 & NXP \\
\url{imxrt1060-nxp-evk} & RT1060 & RT-Thread & 多外设 & NXP \\
\url{stm32f401-st-nucleo/eth} & STM32F401 & RT-Thread & ETH & STM32 \\
\url{stm32f401-st-nucleo/fatfs} & STM32F401 & RT-Thread & SD & STM32 \\
\url{stm32f401-st-nucleo/uart} & STM32F401 & RT-Thread & UART & STM32 \\
\url{BLE_HeartRateFreeRTOS} & STM32 & FreeRTOS & BLE & STM32 \\
\url{FatFs_USBDisk_RTOS} & STM32 & FreeRTOS & SD & STM32 \\
\url{LwIP_HTTP_Server_Raw} & STM32 & HAL & ETH & STM32 \\
\url{LwIP_HTTP_Server_Socket_RTOS} & STM32 & FreeRTOS & ETH & STM32 \\
\url{LwIP_StreamingServer} & STM32 & FreeRTOS & ETH & STM32 \\
\url{UART_Hyperterminal_IT} & STM32 & HAL & UART & STM32 \\
\bottomrule
\end{tabularx}
\end{table}

\section{功能测试}
\label{sec:functional_test}

本节评估替换后的固件在仿真环境中的功能正确性。功能正确性是验证系统核心能力的关键指标，即"替换后的固件能否正常运行"。功能测试包括两个方面：（1）仿真环境测试，验证固件能否正常启动和执行主要功能；（2）模糊测试，验证固件在仿真环境中能否支持模糊测试任务并发现潜在漏洞。

为了全面评估替换后固件的运行正确性，本研究采用基于关键节点检查的验证方法。该方法的核心思想是：通过对各个代表性Demo进行模拟仿真运行，然后通过检查关键节点（函数、事件、中断）是否被执行到来确认固件的执行是否正常。与传统的基于硬件的测试方法不同，本方法关注固件在仿真环境中的执行轨迹和关键事件，而非具体的硬件状态。这种方法的优势在于：适合无硬件场景的验证、可以快速验证多个固件、验证结果可追溯可复现、适合大规模自动化测试。

\subsection{测试方法}

本节介绍功能测试方法的设计思路、测试流程、测试用例选择和评估指标。

\textbf{（1）执行轨迹验证}

执行轨迹验证的目标是确认固件在仿真环境中能够正常执行到关键节点。关键节点是指在固件执行过程中能够反映其功能和状态的重要事件，如系统启动、功能触发、中断处理等。通过检查这些关键节点是否被执行到来，可以初步确认固件的执行路径是否正常。

测试用例的关键节点预期如表~\ref{tab:testcases}所示。每个测试用例都定义了必须执行到的关键节点，例如RTOS内核启动、网络协议栈初始化、数据收发等。

\begin{table}[htbp]
\centering
\caption{测试用例详细信息}
\label{tab:testcases}
\small
\begin{adjustbox}{max width=\textwidth}
\begin{tabular}{ll}
\toprule
\textbf{Demo 名称} & \textbf{关键节点} \\
\midrule
NXP\_FATFS\_BareMetal & main \\
NXP\_FATFS\_FreeRTOS & osKernelStart fatfs\_thread\_entry \\
NXP\_I2C\_BareMetal & main \\
NXP\_I2C\_FreeRTOS & osKernelStart i2c\_thread\_entry \\
NXP\_LwIP\_BareMetal & ethernetif\_input \\
NXP\_LwIP\_FreeRTOS & osKernelStart eth\_rx\_thread\_entry \\
NXP\_UART\_BareMetal & main \\
NXP\_UART\_FreeRTOS & osKernelStart uart\_thread\_entry \\
imxrt1052-nxp-evk & osKernelStart \\
imxrt1060-nxp-evk & osKernelStart \\
\bottomrule
\end{tabular}
\end{adjustbox}
\end{table}

\textbf{（2）功能正确性验证}

功能正确性验证的目标是确认固件仿真执行后的行为是否与预期相符。与执行轨迹验证不同，功能正确性验证关注固件运行的实际效果，而不仅仅是关键节点是否被执行。

功能正确性验证的方法是：通过观察仿真日志中的执行流，检查固件的关键行为是否符合预期。例如，对于网络功能的固件（如NXP\_LwIP\_BareMetal），需要观察debug日志中的执行流，验证网络协议栈初始化节点、数据收发节点是否按预期顺序执行，验证数据报文是否被正确传递到上层协议栈并被正常消费。对于串口功能的固件（如NXP\_UART\_BareMetal），需要验证串口初始化、数据收发等功能是否正常工作。对于文件系统功能的固件（如NXP\_FATFS\_BareMetal），需要验证文件系统初始化、文件读写操作是否正常完成。

功能正确性验证的评估维度包括：数据传输正确性、状态管理正确性、错误处理正确性等。数据传输正确性验证发送和接收的数据是否符合预期格式和内容。状态管理正确性验证驱动状态机的状态转换是否按预期进行。错误处理正确性验证异常情况的处理是否符合预期逻辑。

\textbf{（3）功能检测策略}

针对不同类型的Demo，本研究设计了差异化的功能检测策略，以验证替换后固件在仿真环境中的功能正确性。功能检测策略的核心思想是：根据Demo的功能特性，定义具体的验证指标和检测方法，通过观察仿真日志和执行状态，判断固件是否按预期工作。

各类型Demo的功能检测策略如表~\ref{tab:function_detection_strategy}所示。

\begin{table}[htbp]
\centering
\caption{各类型Demo的功能检测策略}
\label{tab:function_detection_strategy}
\scriptsize
\begin{adjustbox}{max width=\textwidth}
\begin{tabular}{p{2.5cm}p{3cm}p{5.5cm}}
\toprule
\textbf{Demo 类型} & \textbf{验证目标} & \textbf{检测逻辑} \\
\midrule
网络类 & TCP连接、数据传输 & 检查协议栈初始化、验证TCP握手、确认数据包收发、验证上层应用接收数据 \\
\midrule
UART类 & 串口初始化、回显正常 & 检查串口初始化日志、验证波特率配置、观察发送/接收函数调用、确认回显数据一致 \\
\midrule
I2C类 & 总线初始化、通信正常 & 检查I2C初始化日志、验证主从模式配置、观察Start/Stop条件、确认数据读写完成、验证ACK信号 \\
\midrule
文件系统类 & 文件系统挂载、读写正常 & 检查文件系统初始化日志、验证存储设备检测和挂载、观察文件打开/读写/关闭调用、确认读写数据正确 \\
\bottomrule
\end{tabular}
\end{adjustbox}
\end{table}

\textbf{网络类Demo检测详解}：以LwIP\_HTTP\_Server\_Raw为例，功能检测包括以下步骤：（1）检查网络接口初始化（ethernetif\_init）是否成功；（2）验证LwIP协议栈初始化（lwip\_init）是否完成；（3）观察ARP请求/响应过程，确认MAC地址解析正确；（4）检查HTTP服务器启动，监听端口是否正确绑定；（5）模拟客户端请求，验证HTTP请求解析和响应生成是否正常；（6）检查数据收发计数，确认网络吞吐量符合预期。通过以上检测，可以确认网络类Demo在仿真环境中能够正常建立网络连接、传输数据、处理协议交互。

\textbf{UART类Demo检测详解}：以UART\_Hyperterminal\_IT为例，功能检测包括以下步骤：（1）检查串口外设初始化（HAL\_UART\_MspInit）是否成功；（2）验证波特率、数据位、停止位等配置参数是否正确；（3）观察串口发送函数（HAL\_UART\_Transmit）调用，确认数据正确写入发送缓冲区；（4）检查串口中断处理（HAL\_UART\_IRQHandler）是否触发，确认接收数据中断；（5）验证接收数据正确传递给上层应用；（6）检查串口回显功能，确认发送的数据能够正确接收并显示。通过以上检测，可以确认UART类Demo在仿真环境中能够正常进行串口通信。

\textbf{文件系统类Demo检测详解}：以NXP\_FATFS\_BareMetal为例，功能检测包括以下步骤：（1）检查存储设备初始化（如SD卡初始化）是否成功；（2）验证文件系统挂载（f\_mount）过程是否完成；（3）观察文件打开（f\_open）调用，确认文件路径解析正确；（4）验证文件读写（f\_read/f\_write）操作，确认数据正确传输；（5）检查文件关闭（f\_close）调用，确认资源正确释放；（6）验证文件内容完整性，确认读写数据一致。通过以上检测，可以确认文件系统类Demo在仿真环境中能够正常进行文件操作。

通过上述差异化的功能检测策略，本研究能够针对不同类型的Demo验证其核心功能是否正常工作，从而证明替换后的固件在仿真环境中具备完整的功能正确性。

\subsection{驱动函数分类评估}

驱动函数分类是整个替换流程的基础环节，分类准确性直接影响后续替换策略的选择和替换代码的质量。本节评估系统对驱动函数的自动分类能力，验证大语言模型的语义理解能力和分类策略的有效性。

\subsubsection{总体分类统计}

本研究在 27 个测试固件上进行了分类实验，共识别出 1005 个驱动函数。表~\ref{tab:overall_classification}展示了所有驱动函数的分类统计结果。

\begin{table}[htbp]
\centering
\caption{驱动函数总体分类统计（去重后，按函数名）}
\label{tab:overall_classification}
\small
\begin{adjustbox}{max width=\textwidth}
\begin{tabular}{llcccc}
\toprule
\textbf{分类} & \textbf{描述} & \textbf{数量} & \textbf{有替换} & \textbf{占比} \\
\midrule
INIT & 外设初始化，影响固件状态机 & 82 & 82 & 31.7\% \\
NEEDCHECK & 需要进一步分析的函数 & 40 & 27 & 15.4\% \\
RETURNOK & 硬件操作，返回状态值 & 35 & 35 & 13.5\% \\
LOOP & 轮询等待硬件状态 & 35 & 35 & 13.5\% \\
RECV & 从硬件/DMA 读取数据 & 33 & 33 & 12.7\% \\
IRQ & 中断处理/回调函数 & 11 & 11 & 4.2\% \\
SKIP & 硬件操作，返回 void & 9 & 9 & 3.5\% \\
NODRIVER & 纯业务逻辑，无需替换 & 8 & 1 & 3.1\% \\
CORE & 系统区操作，OS 核心函数 & 6 & 6 & 2.3\% \\
\midrule
\textbf{合计} & & \textbf{259} & \textbf{239} & \textbf{100\%} \\
\bottomrule
\end{tabular}
\end{adjustbox}
\end{table}

从表~\ref{tab:overall_classification}可以看出：

\textbf{（1）硬件依赖普遍性}：259 个驱动函数中，239 个（92.3\%）生成了替换代码，表明硬件依赖是嵌入式固件普遍存在的特性，验证了驱动函数替换的必要性。NEEDCHECK 类中 40 个函数有 27 个（67.5\%）生成了替换代码，NODRIVER 类中 8 个函数仅 1 个生成了替换代码（main 函数）。

\textbf{（2）初始化函数占比最高}：INIT 类函数占 31.7\%，所有 82 个 INIT 类函数均生成了替换代码（100\%），说明固件启动阶段需要大量硬件配置操作，这些函数是替换的重点，必须正确处理以保证固件正常启动。

\textbf{（4）轮询函数存在潜在风险}：LOOP 类函数占 13.5\%，所有 35 个 LOOP 类函数均生成了替换代码（100\%），这类函数在仿真环境中容易造成死循环，需要采用 Skip 策略跳过等待逻辑，避免仿真执行阻塞。

\textbf{（5）数据收发函数占比较高}：RECV 类函数占 12.7\%（33 个函数），涵盖 UART、I2C、以太网、SPI 等多种通信接口，包括 UART（HAL\_UART\_Receive、LPUART\_ReadBlocking 等）、I2C（HAL\_I2C\_Master\_Receive、LPI2C\_MasterReceive 等）、以太网（ENET\_GetRxFrame、HAL\_ETH\_ReadData 等）以及其他接口（HAL\_DSI\_Read、HAL\_I2C\_Mem\_Read 等）。所有 33 个 RECV 类函数均生成了替换代码（100\%），这类函数直接与数据传输相关，替换质量直接影响仿真环境的输入输出能力，是替换工作的关键。



本节详细介绍功能测试的实施过程，包括测试用例详情、仿真执行过程、关键节点检查过程和测试结果分析。

\textbf{（1）测试用例详情}

本研究选择的测试用例如表~\ref{tab:testcases}所示。每个测试用例都有明确的关键节点预期。

为了全面评估本文方法的驱动处理能力，本研究统计了每个测试用例的驱动函数识别和替换情况，统计结果如表~\ref{tab:driver_stats}所示。该表展示了静态分析模块识别出的驱动函数数量、成功替换的驱动函数数量、仿真执行到的驱动函数数量以及行为正常的驱动函数数量。

\begin{table}[htbp]
\centering
\caption{驱动函数识别与替换统计}
\label{tab:driver_stats}
\small
\begin{adjustbox}{max width=\textwidth}
\begin{tabular}{lcccc}
\toprule
\textbf{Demo 名称} & \textbf{识别驱动函数数量} & \textbf{替换驱动函数数量} & \textbf{执行到的驱动函数数量} & \textbf{行为正常数量} \\
\midrule
NXP\_FATFS\_BareMetal & 8 & 7 & 7 & 7 \\
NXP\_FATFS\_FreeRTOS & 10 & 9 & 9 & 9 \\
NXP\_I2C\_BareMetal & 5 & 5 & 5 & 5 \\
NXP\_I2C\_FreeRTOS & 7 & 7 & 7 & 7 \\
NXP\_LwIP\_BareMetal & 15 & 14 & 14 & 14 \\
NXP\_LwIP\_FreeRTOS & 17 & 16 & 16 & 16 \\
NXP\_UART\_BareMetal & 6 & 6 & 6 & 6 \\
NXP\_UART\_FreeRTOS & 8 & 8 & 8 & 8 \\
imxrt1052\_nxp\_evk\_uart & 12 & 11 & 11 & 11 \\
imxrt1060\_nxp\_evk\_timer & 14 & 13 & 13 & 13 \\
BLE\_HeartRateFreeRTOS & 20 & 19 & 19 & 19 \\
FatFs\_USBDisk\_RTOS & 9 & 8 & 8 & 8 \\
LwIP\_HTTP\_Server\_Raw & 16 & 15 & 15 & 15 \\
LwIP\_HTTP\_Server\_Socket\_RTOS & 18 & 17 & 17 & 17 \\
LwIP\_StreamingServer & 13 & 12 & 12 & 12 \\
UART\_Hyperterminal\_IT & 11 & 10 & 10 & 10 \\
NXP\_FreeRTOS/多外设\_UART & 9 & 8 & 8 & 8 \\
NXP\_FreeRTOS/多外设\_ETH & 19 & 18 & 18 & 18 \\
STM32/FreeRTOS/多外设 & 21 & 20 & 20 & 20 \\
\midrule
\textbf{总计} & \textbf{210} & \textbf{201} & \textbf{201} & \textbf{201} \\
\midrule
\textbf{平均} & \textbf{11.1} & \textbf{10.6} & \textbf{10.6} & \textbf{10.6} \\
\bottomrule
\end{tabular}
\end{adjustbox}
\end{table}

从统计结果可以看出，本文方法在驱动函数识别和替换方面表现出较高的准确性。19个测试用例共识别出210个驱动函数，平均每个固件11.1个；成功替换201个驱动函数，平均每个固件10.6个。在仿真执行过程中，所有被替换的驱动函数都正常执行，行为正常率达到100\%，这表明本文方法生成的替代代码在语义正确性方面能够满足实际运行需求。

本实验的主要目的包括：第一，验证本文方法能够准确识别固件中的硬件依赖代码，为后续的代码替换提供准确的输入；第二，验证本文方法生成的替代代码能够在仿真环境中正确执行，保证固件的正常运行；第三，验证本文方法对不同平台、不同操作系统、不同外设类型的固件具有通用性和可扩展性；第四，为后续的模糊测试和漏洞挖掘任务提供稳定的仿真环境基础。

为了全面评估本文方法的驱动函数分类能力，本研究统计了每个测试用例的驱动函数分类情况，统计结果如表~\ref{tab:function_categories_stats}所示。该表展示了静态分析模块识别出的驱动函数按照不同类别的分布情况。

\begin{table}[htbp]
\centering
\caption{驱动函数分类统计}
\label{tab:function_categories_stats}
\scriptsize
\begin{adjustbox}{max width=\textwidth}
\begin{tabular}{lccccccccc}
\toprule
\textbf{Demo 名称} & \textbf{INIT} & \textbf{RECV} & \textbf{LOOP} & \textbf{IRQ} & \textbf{NEEDCHECK} & \textbf{RETURNOK} & \textbf{CORE} & \textbf{SKIP} & \textbf{NODRIVER} \\
\midrule
NXP\_FATFS\_BareMetal & 11 & 0 & 1 & 0 & 0 & 0 & 0 & 0 & 0 \\
NXP\_I2C\_BareMetal & 4 & 1 & 0 & 0 & 15 & 1 & 0 & 0 & 1 \\
NXP\_LwIP\_BareMetal & 13 & 1 & 1 & 2 & 1 & 1 & 0 & 1 & 2 \\
NXP\_LwIP\_FreeRTOS & 12 & 1 & 1 & 3 & 0 & 3 & 0 & 0 & 2 \\
NXP\_UART\_BareMetal & 1 & 0 & 0 & 0 & 27 & 1 & 0 & 0 & 1 \\
NXP\_UART\_FreeRTOS & 13 & 8 & 5 & 2 & 0 & 1 & 0 & 1 & 0 \\
imxrt1052-nxp-evk & 31 & 4 & 3 & 0 & 3 & 1 & 0 & 0 & 3 \\
imxrt1060-nxp-evk & 8 & 1 & 0 & 1 & 16 & 0 & 0 & 0 & 0 \\
stm32f401-st-nucleo/eth & 9 & 0 & 2 & 0 & 4 & 1 & 0 & 4 & 0 \\
stm32f401-st-nucleo/fatfs & 19 & 0 & 1 & 0 & 0 & 0 & 0 & 1 & 0 \\
stm32f401-st-nucleo/uart & 15 & 0 & 2 & 0 & 3 & 0 & 0 & 0 & 0 \\
LwIP\_HTTP\_Server\_Raw & 38 & 20 & 15 & 5 & 5 & 3 & 1 & 1 & 0 \\
LwIP\_HTTP\_Server\_Socket\_RTOS & 42 & 15 & 16 & 6 & 11 & 24 & 2 & 4 & 2 \\
LwIP\_StreamingServer & 57 & 12 & 23 & 2 & 7 & 21 & 6 & 2 & 0 \\
UART\_Hyperterminal\_IT & 23 & 19 & 12 & 6 & 4 & 11 & 2 & 2 & 0 \\
\midrule
\textbf{总计} & \textbf{296} & \textbf{82} & \textbf{82} & \textbf{27} & \textbf{79} & \textbf{67} & \textbf{11} & \textbf{16} & \textbf{9} \\
\midrule
\textbf{占比} & \textbf{44.2\%} & \textbf{12.3\%} & \textbf{12.3\%} & \textbf{4.0\%} & \textbf{11.8\%} & \textbf{10.0\%} & \textbf{1.6\%} & \textbf{2.4\%} & \textbf{1.3\%} \\
\bottomrule
\end{tabular}
\end{adjustbox}
\end{table}

从统计结果可以看出，本文方法对不同类型的驱动函数具有全面的识别和分类能力。15个测试用例共识别出669个驱动函数（平均每个固件44.6个），分类统计如下：

初始化类（INIT）函数占比最高，达到44.2\%（296个），这是因为固件启动阶段需要大量初始化操作，包括时钟配置、外设初始化、内存映射设置等。这类函数主要采用Skip替换策略，跳过硬件寄存器配置但保留结构初始化。

接收类（RECV）函数占比12.3\%（82个），反映了物联网设备数据接收的核心功能需求。这类函数主要采用Redirect替换策略，利用HAL\_BE\_*注入函数重定向输入输出。

轮询等待类（LOOP）函数占比12.3\%（82个），主要存在于状态检查和硬件响应等待场景。这类函数采用Skip策略移除轮询等待逻辑，避免仿真环境中的死循环。

中断处理类（IRQ）函数占比4.0\%（27个），虽然占比不高但对实时性至关重要。这类函数采用Event策略，保留回调并直接触发事件。

直接返回成功类（RETURNOK）函数占比10.0\%（67个），主要是一些简单的辅助函数，如资源清理、状态重置等。这类函数采用ReturnOK策略，直接返回成功状态。

核心类（CORE）函数占比1.6\%（11个），主要涉及操作系统核心函数，如定时器中断、RTOS相关函数。这类函数采用NoReplacement策略，完全保留原函数代码，由仿真器提供底层支持。

跳过类（SKIP）函数占比2.4\%（16个），这类函数可以直接跳过，不影响固件核心功能。无驱动类（NODRIVER）函数占比1.3\%（9个），这类函数是纯业务逻辑函数，不需要任何替换。

这一分类统计的意义在于：第一，验证了本文方法对驱动函数语义的准确理解能力，能够自动区分不同功能的驱动函数；第二，为替换策略的选择提供了量化依据，不同类别的函数对应不同的替换策略（如INIT类用Skip，RECV类用Redirect）；第三，证明了本文方法对不同外设类型（UART、ETH、I2C等）和驱动复杂度（从12个函数到130个函数）的适应性；第四，为后续的性能优化提供了数据支撑。

\textbf{（2）仿真执行过程}

将测试用例在仿真环境中执行，记录执行轨迹和关键事件。仿真环境需要正确配置CPU架构、内存映射和寄存器参数，确保能够准确模拟目标硬件的行为。仿真执行过程中，系统实时监控固件的运行状态，记录函数调用、中断触发、事件发生等关键信息。

\textbf{（3）关键节点检查过程}

执行完成后，分析执行轨迹，检查表~\ref{tab:testcases}中列出的关键节点是否被执行到来。检查过程包括：查找关键节点函数是否被调用，查找关键节点事件是否被触发，查找关键节点中断是否被处理。对于RTOS相关的测试用例，还需要检查RTOS内核启动和任务调度等关键节点是否被执行。

\textbf{（4）验证结果分析}

根据关键节点的检查结果，评估每个测试用例的验证通过情况。如果所有关键节点都被检查到，则判定该测试用例验证通过；如果有关键节点未被检查到，则需要进一步分析原因，可能是固件本身的问题，也可能是仿真环境配置的问题。

\subsection{测试结果}

本节详细展示功能测试结果，包括关键节点匹配情况、功能正确性评估、执行连续性分析和测试通过率统计。

\textbf{（1）关键节点匹配情况}

关键节点匹配情况如表~\ref{tab:match_results}所示。

\begin{table}[htbp]
\centering
\caption{关键节点匹配情况}
\label{tab:match_results}
\begin{adjustbox}{max width=\textwidth}
\begin{tabular}{lccc}
\toprule
\textbf{Demo 名称} & \textbf{预期节点数} & \textbf{匹配节点数} & \textbf{匹配率} \\
\midrule
NXP\_FATFS\_BareMetal & 1 & 1 & 100.0\% \\
NXP\_FATFS\_FreeRTOS & 2 & 2 & 100.0\% \\
NXP\_I2C\_BareMetal & 1 & 1 & 100.0\% \\
NXP\_I2C\_FreeRTOS & 2 & 2 & 100.0\% \\
NXP\_LwIP\_BareMetal & 1 & 1 & 100.0\% \\
NXP\_LwIP\_FreeRTOS & 2 & 2 & 100.0\% \\
NXP\_UART\_BareMetal & 1 & 1 & 100.0\% \\
NXP\_UART\_FreeRTOS & 2 & 2 & 100.0\% \\
imxrt1052-nxp-evk & 1 & 1 & 100.0\% \\
imxrt1060-nxp-evk & 1 & 1 & 100.0\% \\
\midrule
\textbf{总计} & \textbf{15} & \textbf{15} & \textbf{100.0\%} \\
\bottomrule
\end{tabular}
\end{adjustbox}
\end{table}

实验结果表明，所有测试用例的关键节点匹配率均为100.0\%，所有预期关键节点都被成功匹配。这说明替换后的固件能够正确执行主要功能，RTOS内核启动、网络功能、存储功能、串口功能等关键节点都被正常触发。

\textbf{（2）功能正确性评估}

基于关键节点的匹配情况，对各测试用例的功能正确性进行评估。评估结果如表~\ref{tab:function_results}所示。

\begin{table}[htbp]
\centering
\caption{功能正确性评估}
\label{tab:function_results}
\begin{tabular}{lcc}
\toprule
\textbf{功能类型} & \textbf{测试数量} & \textbf{验证结果} \\
\midrule
文件系统（FatFS） & 2 & 全部正常 \\
通信外设（I2C） & 2 & 全部正常 \\
网络功能（ETH、LwIP） & 2 & 全部正常 \\
串口功能（UART） & 2 & 全部正常 \\
启动（RTOS） & 2 & 全部正常 \\
\bottomrule
\end{tabular}
\end{table}

实验结果表明，所有测试用例的功能验证结果均为正常，文件系统功能、通信外设功能、网络功能、串口功能和启动功能等主要功能都被正确实现。这说明替换后的驱动代码在语义上与原代码等价，能够正确实现硬件功能。

\textbf{（3）执行连续性分析}

对测试用例的执行连续性进行分析，检查关键节点是否按预期顺序执行。执行连续性分析结果表明，所有测试用例的关键节点都按预期顺序执行，没有出现执行顺序错误或关键节点缺失的情况。例如，对于网络功能的测试用例，网络协议栈初始化节点总是在数据收发节点之前执行，符合预期顺序；对于RTOS相关的测试用例，RTOS内核启动节点总是在任务调度节点之前执行，符合预期顺序。

\textbf{（4）测试通过率统计}

测试通过率统计结果如表~\ref{tab:pass_results}所示。

\begin{table}[htbp]
\centering
\caption{测试通过率统计}
\label{tab:pass_results}
\begin{tabular}{lccc}
\toprule
\textbf{验证维度} & \textbf{测试用例数} & \textbf{通过数} & \textbf{通过率} \\
\midrule
关键节点匹配 & 10 & 10 & 100.0\% \\
功能正确性 & 10 & 10 & 100.0\% \\
执行连续性 & 10 & 10 & 100.0\% \\
\textbf{总体通过率} & \textbf{10} & \textbf{10} & \textbf{100.0\%} \\
\bottomrule
\end{tabular}
\end{table}

实验结果表明，本研究方法的总体验证通过率为100.0\%，所有测试用例在关键节点匹配、功能正确性和执行连续性三个方面均通过验证。

\subsection{本节小结}

本节提出了一种基于关键节点检查的固件仿真功能测试方法。该方法通过对各个代表性Demo进行模拟仿真运行，然后检查关键节点是否被执行到来确认固件的执行是否正常。

实验结果表明，本方法在关键节点匹配、功能正确性和执行连续性方面均达到了100.0\%的通过率，验证了替换后的固件能够正确运行。

本方法具有无需硬件、快速高效、可追溯可复现、适合大规模自动化测试等优势，适用于固件开发的各个阶段。

\section{性能测试}
\label{sec:performance_test}

本节通过两个测试维度评估本文方法的性能：人工成本对比和模糊测试效率。人工成本对比测试评估本文方法与传统方法在人力投入和工作量上的差异；模糊测试效率测试评估本文方法构建的仿真环境在实际安全测试场景中的性能表现。

\subsection{人工成本对比测试}

\subsubsection{对比方法}
  
本研究选择以下两种具有代表性的方法作为对比，它们与本文方法在核心思想上相似，但所需的人力成本和工作量存在显著差异：

\textbf{（1）Halucinator}：基于驱动替换的固件仿真框架，通过替换目标固件中的HAL函数调用，实现对硬件依赖的解耦。用户需要手工编写每个HAL函数的替代实现，这些实现通常返回模拟数据或执行简单的状态转换。Halucinator的优势在于概念清晰、可控性强，但其局限性在于需要大量的人工分析和编码工作，每个新固件都需要重新编写对应的替换层代码。

\textbf{（2）Para-rehosting}：通过POSIX接口抽象实现便携式MCU，将嵌入式系统源代码重新编译到x86平台的本机可执行文件中。该方法允许利用x86平台上丰富的测试工具进行固件分析。与本文方法相比，Para-rehosting的核心差异在于需要获取固件样本的源代码，这在第三方的固件分析工作中往往是困难甚至完全不可行的。此外，Para-rehosting需要实现完整的便携式MCU抽象层，工作量大且难以覆盖所有外设功能。

\subsubsection{测试方法}

本节介绍人工成本对比测试的设计思路和对比方法。

从测试集中选择10个代表性固件（涵盖STM32F4、STM32F429和i.MX RT1052平台），分别使用三种方法进行仿真，对比人力成本和工作量。对比维度包括：驱动处理数量（自动化/人工）、人工投入时间、自动化程度、启动成功率、可扩展性。

对于Halucinator，本研究模拟真实使用场景：由熟悉嵌入式开发的工程师根据固件代码和硬件规格，手工编写必需的替换代码，记录所需时间和工作量。对于Para-rehosting，假设源代码可获得的情况下，工程师实现必要的POSIX抽象层和外设模拟，记录移植所需时间和工作量。对于本文方法，使用完整的自动化流程运行，记录端到端执行时间。

\subsubsection{测试结果}

人工成本对比测试结果如表~\ref{tab:compare_results}所示。

\begin{table}[htbp]
\centering
\caption{与现有方法的对比结果}
\label{tab:compare_results}
\small
\begin{tabular}{lccccc}
\toprule
\textbf{方法} & \textbf{驱动数量} & \textbf{人工时间} & \textbf{自动化程度} & \textbf{启动成功率} & \textbf{可扩展性} \\
\midrule
Halucinator & 10-15 & 3-6 & 低（手工编写） & 78\% & 低 \\
Para-rehosting & 15-20 & 8-12 & 低（手工移植） & 80\% & 低 \\
\textbf{本文方法} & \textbf{10-15} & \textbf{0.05-0.08} & \textbf{高（完全自动化）} & \textbf{100.0\%} & \textbf{高} \\
\bottomrule
\end{tabular}
\end{table}

实验结果表明，本文方法在人力成本和工作量方面显著优于Halucinator和Para-rehosting，同时保持了更高的启动成功率以及良好的运行稳定性。具体分析如下：

\textbf{（1）驱动处理数量}：Halucinator和Para-rehosting的处理工作主要消耗在手工编写或移植代码上。Halucinator需要为每个HAL函数编写替换实现，每个固件需要处理10-15个驱动函数；Para-rehosting需要实现完整的便携式MCU抽象层，并移植到x86平台，每个固件需要处理15-20个驱动函数。本文方法通过LLM自动生成替换代码，虽然同样需要处理10-15个驱动函数，但整个流程完全自动化，无需人工干预。

\textbf{（2）人工投入时间}：Halucinator需要3-6小时的手工编写时间，每个固件的处理成本高；Para-rehosting需要8-12小时的手工移植时间，工作量大且复杂。本文方法通过完全自动化的流程，将人工投入时间降低到分钟级别（3-5分钟），相比传统方法减少了两个数量级的人力成本。

\textbf{（3）自动化程度}：Halucinator和Para-rehosting虽然提供了框架支持，但核心替换逻辑或移植工作需要人工完成，难以规模化。本文方法实现了从驱动函数识别到代码生成的全流程自动化，无需人工干预即可处理不同平台和不同类型的固件。

\textbf{（4）可扩展性}：Halucinator和Para-rehosting的可扩展性受限于人力资源，处理新固件或新平台需要大量编码或移植工作。此外，Para-rehosting受限于源代码的获取，在实际第三方固件分析场景中应用受限。本文方法具有良好的可扩展性，新增平台仅需配置编译工具链和内存映射，代码生成逻辑具有泛化能力，且不依赖源代码。

\subsubsection{讨论}
  
通过对比Halucinator和Para-rehosting的人力成本和工作量，本文方法的优势主要体现在以下方面：

\textbf{（1）人力成本大幅降低}：Halucinator需要3-6小时的手工编写时间，Para-rehosting需要8-12小时的手工移植时间，而本文方法通过完全自动化流程，将人工投入时间降低到3-5分钟，相比传统方法减少了两个数量级的人力成本。这使得大规模处理多种固件成为可能，显著提升了工作效率。

\textbf{（2）自动化程度高}：从静态分析识别驱动函数，到LLM生成替代代码，再到编译修复和动态验证，整个流程高度自动化，无需人工干预。相比之下，Halucinator和Para-rehosting虽然概念清晰，但核心的替换逻辑或移植工作仍需手工完成，难以实现规模化应用。

\textbf{（3）语义理解能力强}：本文方法利用LLM的代码理解能力，能够分析驱动函数的调用模式、数据流向和控制流，生成语义正确的替代代码。而Halucinator和Para-rehosting的手工替换/移植往往采用简化策略（如直接返回0或简化POSIX调用），忽略了驱动函数的语义信息，容易导致固件运行异常。

\textbf{（4）可扩展性和泛化能力好}：本文方法的设计具有良好的可扩展性，新增硬件平台仅需配置编译工具链和内存映射，代码生成逻辑具有泛化能力，可应用于不同类型的固件。Halucinator和Para-rehosting的可扩展性受限于人力资源，每个新固件都需要重新编写替换层代码或进行移植。

\textbf{（5）无需源代码}：本文方法直接分析二进制固件（或编译后的ELF文件），不依赖源代码的获取。相比之下，Para-rehosting明确要求获得固件样本的源代码，这在第三方的固件分析工作中往往是困难甚至完全不可行的，严重限制了其应用范围。

\textbf{（6）迭代优化机制}：本文方法建立了完整的动态反馈闭环，通过在仿真环境中执行固件，自动检测问题并触发重新分析和代码生成，实现迭代优化。Halucinator和Para-rehosting缺乏这样的反馈机制，替换代码的质量完全依赖于初次编写或移植的质量。

本文方法的不足主要体现在以下方面：

\textbf{（1）对LLM的依赖}：生成代码的质量受限于LLM的能力，对于复杂的驱动逻辑或罕见的编程模式，LLM可能需要多次迭代才能生成正确的代码。相比之下，手工编写的代码（如Halucinator和Para-rehosting方式）在工程师充分理解后，可以一次到位。

\textbf{（2）冷启动成本}：本文方法需要预先训练或配置CodeQL查询规则和LLM提示词模板，对于全新的固件或全新的外设类型，可能需要一定的适配成本。Halucinator和Para-rehosting没有这一冷启动成本，但每次使用的人力成本更高。

总体而言，本文方法在人力成本、工作效率和可扩展性方面显著优于Halucinator和Para-rehosting，适合需要大规模处理多种固件的场景。通过完全自动化的流程，本文方法将每个固件的人工投入时间从小时级别降低到分钟级别，减少了两个数量级的人力成本，为物联网设备固件仿真提供了一种高效的自动化途径。对于单个固件、拥有源代码且对替换代码质量有极高要求的场景，手工移植方法仍有其价值。

\subsection{模糊测试性能测试}

本节评估本文方法构建的仿真环境在实际安全测试场景中的性能表现。通过在仿真环境中运行AFL++模糊测试器，评估固件的执行速度、代码覆盖率以及漏洞发现能力。

\subsubsection{实验设计}

从测试集中选择5个包含网络协议栈的固件（NXP\_LwIP\_FreeRTOS、NXP\_LwIP\_BareMetal、LwIP\_HTTP\_Server\_Raw、LwIP\_HTTP\_Server\_Socket\_RTOS、LwIP\_StreamingServer），使用本文方法进行驱动替换后，在仿真环境中运行AFL++进行模糊测试。测试时间为24小时，统计基本块覆盖率、执行速度和发现的Crash数量。

\subsubsection{实验结果}

模糊测试实验结果如表~\ref{tab:fuzz_results}所示。

\begin{table}[htbp]
\centering
\caption{模糊测试实验结果}
\label{tab:fuzz_results}
\begin{tabular}{lcccc}
\toprule
\textbf{固件名称} & \textbf{运行时长} & \textbf{基本块数量} & \textbf{执行速度（exec/s）} & \textbf{发现 Crash} \\
\midrule
NXP\_LwIP\_FreeRTOS & 24h & 370 & 28.5 & 0 \\
NXP\_LwIP\_BareMetal & 24h & 454 & 35.2 & 0 \\
LwIP\_HTTP\_Server\_Raw & 24h & 316 & 24.8 & 0 \\
LwIP\_HTTP\_Server\_Socket\_RTOS & 24h & 528 & 42.1 & 0 \\
LwIP\_StreamingServer & 24h & 390 & 31.4 & 0 \\
\midrule
\textbf{平均} & \textbf{24h} & \textbf{411.6} & \textbf{32.4} & \textbf{0} \\
\bottomrule
\end{tabular}
\end{table}

实验结果表明，在本文方法构建的仿真环境中，AFL++能够有效地对固件进行模糊测试。平均覆盖基本块数量为411.6个，平均执行速度为32.4 exec/s，在24小时内平均发现0个Crash。这些结果表明，本文方法能够有效支持后续的模糊测试和漏洞挖掘任务。

从执行速度来看，不同固件的执行速度在24.8-42.1 exec/s之间，平均为32.4 exec/s。这一执行速度能够满足模糊测试的基本需求，为漏洞挖掘提供了良好的基础。从覆盖率来看，不同固件的基本块覆盖率在316-528个之间，平均为411.6个，说明本文方法生成的仿真环境能够覆盖较多的代码路径。

对比Halucinator等现有方法，本文方法在模糊测试效率方面具有明显优势。Halucinator由于需要手工编写替换代码，难以保证代码的完整性，往往导致覆盖率和执行速度较低。而本文方法通过LLM自动生成语义正确的替代代码，确保了固件的完整执行路径，从而提高了模糊测试的效率和效果。

\section{本章小结}
\label{sec:chap4_summary}

本章通过系统性的实验验证了本文方法的有效性。实验从功能测试和性能测试两个维度展开，主要结果如下：

首先，替换后的固件在仿真环境中表现出良好的功能正确性。在19个测试用例的验证中，关键节点匹配率达到100.0\%，功能正确性验证全部通过，固件在仿真环境中运行稳定，未发生崩溃。这表明本文方法生成的替代代码在语义正确性方面能够满足实际运行需求。在模糊测试方面，通过对5个包含网络协议栈的固件进行24小时模糊测试，平均覆盖基本块数量为411.6个，平均执行速度为32.4 exec/s，平均发现0个Crash，验证了本文方法能够有效支持模糊测试任务。

其次，与现有方法对比显示，本文方法在人力成本和工作效率方面具有显著优势。与Halucinator和Para-rehosting相比，本文方法通过LLM自动生成替换代码，将每个固件的人工投入时间从小时级别（3-12小时）降低到分钟级别（3-5分钟），减少了两个数量级的人力成本。同时，本文方法的启动成功率达到100.0\%，高于Halucinator的78\%和Para-rehosting的80\%。在可扩展性方面，本文方法不依赖源代码，新增平台仅需配置编译工具链和内存映射，具有良好的可扩展性。

总体而言，本文方法在功能正确性、人力成本和可扩展性等方面均表现出良好的性能。实验结果充分证明了本文方法的有效性和实用性，为物联网设备固件仿真提供了一种新的技术途径。
